\section{Problems Found and Pitfalls of Browser Benchmarking}
\label{sec:problems}

A browser is not a process that waits to be measured: it moves itself,
starts work of its own and updates itself. This section lists every problem
the study met, how it was detected and how it is controlled. Several were
found on earlier builds and fixed before the reported one; they are part of
the method and of the contribution.

\begin{table}[htbp]
\caption{The problems found, in the order they appeared.}
\label{tab:problems}
\centering\scriptsize
\begin{tabular}{rp{0.25\linewidth}p{0.3\linewidth}p{0.32\linewidth}}
\toprule
\# & Problem & How it was detected & Control \\
\midrule
1 & A mask of eight CPUs that is four cores & Review of the CPU numbering before measuring & SMT off, one CPU per core, checked per run \\
2 & Background work on the measured cores & Design review & Shield, quiet mode, foreign CPU time per run \\
3 & Asymmetric background services between the browsers & Review of the launch settings & The same services off in both \\
4 & Chrome's browser process leaves the CPU mask & A sampler of every thread's allowed CPUs & No session bus for the browser; escaped processes counted per run and rejected \\
5 & Firefox still starting up during the measurement & Bursty iterations; a pilot with and without delay & A 3 s start delay in every browser run \\
6 & A slow first browser launch & Launch times on a new machine & One discarded launch per browser before the plan \\
7 & A command run during a campaign & One run over the foreign CPU threshold & Nothing else is started during a campaign \\
8 & Our own asymmetries in RQ2 & Review of the measured path against its machine code & A new build; everything measured again \\
9 & The pinned Firefox updated itself & Strict validation of the recorded version & An enterprise policy, the version read from the install, a check before every campaign \\
10 & A degraded core & Crashes pinned to one core & Measure only on the prepared machine; per-core check; the build made twice \\
11 & A wrong computation without a crash & A type error inside a library on valid data & Analysis pinned away from the core; tables regenerated on another machine \\
12 & A new workload changed another's machine code & The codegen comparison & The sequential filter has its own copy of the row code \\
13 & A gate criterion failed twice & The gate & A recorded decision before the campaign; the effect measured \\
14 & Whole-run counters are mostly not the loop & A sanity check against the listing & Counters by the difference of two pass counts \\
\bottomrule
\end{tabular}
\end{table}

\paragraph{4. Chrome escapes the mask.} During the mask pilot a sampler read
the allowed CPUs of every thread and showed Chrome's browser process on the
E-cores. Chrome asks the desktop session, through the D-Bus session bus, to
move it into its own systemd scope, which is outside the run's scope and
its mask. Firefox's processes all stayed inside, so the two browsers were
not confined alike. Without the session bus address every Chrome process
stays in the run scope. Each run now counts the processes of its browser
found outside its scope, and strict validation rejects a run with any. The
renderer, where the WASM and its workers run, had stayed on the P-cores
throughout, and the pilots taken before the fix were repeated where their
decision could depend on it. A process-level mask is a request that a
process can undo.

\paragraph{5. Firefox's start-up work.} With a fresh profile on every run,
the page starts measuring while the browser is still starting. For short
iterations at eight threads this made Firefox's local histogram 1.32 times
slower and bursty, and produced a ``knee'' at seven and eight threads that
looked like oversubscription and was first read as such. A 3~s delay before
the page creates its worker removes both. A scaling artefact at high $N$
that comes from start-up is the kind of effect that would be reported as a
result of the engine.

\paragraph{7. Activity during a campaign.} On the second build one run
exceeded the foreign-CPU threshold because a version-control command ran
during a campaign; the run was repeated. Since then campaigns start quiet
mode themselves. For the v3 campaigns the coding assistant used during
development was kept alive, confined to the E-cores, instead of being
stopped (stopping it also froze the terminals it ran in for the 14 hours of
the main campaign), under the rule that no command is started while a
campaign runs. A rule cannot be checked, but what each run records can: the
load average and the foreign CPU time of the v3 main campaign
(Section~\ref{sec:quality}) are those of the v2 main campaign, in which the
assistant was stopped (medians 3.00 and 3.19, maxima of foreign CPU time
0.205 and 0.222~s). That is evidence of the same quiet, not proof that
nothing ran.

\paragraph{8. Our own asymmetries.} The RQ2 numerator paid a bounds check
per pixel that the denominator did not; and the final step of a histogram
was timed in some variants only. Neither would show in any quality check,
since every checksum was right. They were found by reading the measured
path against its disassembly. The fix was not a correction of the numbers
but a new build and a second full measurement
(Section~\ref{sec:builds-history}).

\paragraph{9. Firefox updates itself.} The profile setting that turns
automatic updates off does not stop the download or the update on the next
start. During the exploratory campaigns of the second build the pinned
Firefox 158.0 replaced itself with 160.0. Because every run records the
version of the binary it launched, strict validation rejected all 2640
affected runs; they were set aside and repeated on 158.0. Only an
enterprise policy file stops the updater; the installer now writes it, reads
the version from the install itself, and a campaign refuses to start with
another version. In v3 no campaign was affected and every environment is
interleaved in every campaign. Chrome for Testing has no updater.

\paragraph{10 and 11. The degraded core.} Section~\ref{sec:core}. The lesson
beyond this machine: a crash is the easy failure. A core that computes
wrongly without crashing is caught only by checksums (in the measured
code), by reproducible builds (in the compiler) and by regenerating results
elsewhere (in the analysis).

\paragraph{12. Code that did not change, changed.} Adding a second caller
to a function that the compiler had inlined changed the machine code of the
parallel filter, whose source was untouched (Section~\ref{sec:codegen}). A
comparison of two builds is a comparison of two instruction streams;
without a tool that compares them, this would have been measured as noise
or, worse, as a result.

\paragraph{14. Counters of a whole run.} Section~\ref{sec:why-rq1}: per-pixel
counters of a browser run were three to eight times what the loop executes,
because a run is also a browser start, an input generation and checksums.
Only the difference between two campaigns that differ in the amount of loop
work isolates the loop.

\paragraph{For benchmarkers.} A CPU mask, a fresh profile and a pinned
version do not by themselves confine, quiet and freeze a browser. What made
these problems visible was recording the state with every run (the allowed
CPUs, the version, the foreign CPU time, the checksum) and rejecting runs
by rules fixed in advance.
